% GENERATED FILE. Edit canonical lessons, then run npm run generate:books.
% canonical-source-hash: fnv1a64:ccfa11bc97114fc4
% canonical-lessons: PT-C13-preto-branco, PT-C13-vermelho-azul

\chapter{Colours}
\label{ch:pt-colours}
\hlchaptermodality{13}

\begin{chapteropening}
\textbf{By the end of this chapter:} \emph{I can name black, white, red and blue in Portuguese and say why not one of the four is a plain inherited Latin colour word.}

Say preto, branco, vermelho and azul, take vermelho back to vermiculus, the little worm crushed for dye, and azul back to Arabic lāzaward.
\end{chapteropening}

\section[preto, branco]{preto, branco — two blacks and a borrowed white}

\label{lesson:PT-C13-preto-branco}



\begin{quote}
Portuguese does something the other Romance languages don't: it keeps \textbf{two words for black}, and the everyday one isn't the Latin one at all.
\end{quote}

\begin{cousinweb}
\noindent
\begin{tabularx}{\linewidth}{@{}>{\raggedright\arraybackslash}X>{\raggedright\arraybackslash}X>{\raggedright\arraybackslash}X@{}}
\toprule
\textbf{word} & \textbf{source} & \textbf{use} \\
\midrule
\textbf{negro} & Latin \emph{\textbf{niger}} & literary, formal, fixed phrases \\
\textbf{preto} & Latin \emph{\textbf{pressus}} \textquotedblleft{}pressed, dense\textquotedblright{} & the \textbf{everyday} colour word \\
\bottomrule
\end{tabularx}

\emph{Negro} is the expected inheritance — the same \emph{niger} that gave French \emph{noir} and Italian \emph{nero}. But day to day, Portuguese reaches for \textbf{preto}, from \emph{pressus} (\textquotedblleft{}pressed together, compact\textquotedblright{}). Dense $\to$ thick $\to$ \textbf{dark}: a colour named for how \textbf{packed} it is, not for its shade.

You already know \emph{pressus}'s English children: \emph{press}, \emph{pressure}, \emph{compress}.
\end{cousinweb}

\begin{cousinweb}
Latin's white, \emph{\textbf{albus}}, lost here too. Portuguese says \textbf{branco}, from Germanic \emph{\textbf{blank}} (\textquotedblleft{}shining\textquotedblright{}) — the same loan that gave French \emph{blanc} and Italian \emph{bianco}, arriving with the Suevi and Visigoths who settled Iberia.

Notice what Portuguese did to the shape: \emph{blanc-} $\to$ \textbf{branc-}. The \textbf{l $\to$ r} swap is a Portuguese signature (compare \emph{plaza} $\to$ \textbf{praça}, \emph{blando} $\to$ \emph{brando}).
\end{cousinweb}

\begin{cousinweb}
\noindent
\begin{tabularx}{\linewidth}{@{}>{\raggedright\arraybackslash}X>{\raggedright\arraybackslash}X@{}}
\toprule
\textbf{word} & \textbf{meaning} \\
\midrule
\textbf{alvorada} & dawn (the sky going white) \\
\textbf{alva} & white, dawn (poetic) \\
\emph{albumina}, \emph{álbum} & (technical borrowings) \\
\bottomrule
\end{tabularx}
\end{cousinweb}

\subsection*{Your turn}
Say these aloud:

\begin{itemize}
\raggedright
  \item \textquotedblleft{}preto, branco\textquotedblright{}
  \item \textquotedblleft{}o vinho branco\textquotedblright{} — Chapter 11's wine
  \item \textquotedblleft{}o pão é branco\textquotedblright{} — Chapter 11's bread
  \item the l$\to$r swap — \textquotedblleft{}blanc $\to$ \textbf{branco}\textquotedblright{}
\end{itemize}

\begin{tcolorbox}[breakable,colback=teal!4,colframe=teal!35!black,title={Before you move on}]
What are Portuguese's two blacks? (\textbf{Negro} $\leftarrow$ \emph{niger}, literary; and \textbf{preto} $\leftarrow$ \emph{pressus} \textquotedblleft{}pressed/dense,\textquotedblright{} the everyday one.) Why would \textquotedblleft{}pressed\textquotedblright{} mean black? (\textbf{Dense $\to$ thick $\to$ dark}.) Where does \textbf{branco} come from? (\textbf{Germanic \emph{blank}}, \textquotedblleft{}shining\textquotedblright{} — borrowed \textbf{into} Romance.) What sound change shows in \emph{branco}? (\textbf{l $\to$ r}, as in \emph{praça}.) Next: \textbf{vermelho}, the strangest colour word in the language.
\end{tcolorbox}
\section[vermelho, azul]{vermelho, azul — a worm and a stone}

\label{lesson:PT-C13-vermelho-azul}



\begin{quote}
These two are the best etymologies in the whole chapter. Portuguese named its red after an \textbf{insect} and its blue after a \textbf{rock}.
\end{quote}

\begin{cousinweb}
Every other Romance language builds \textquotedblleft{}red\textquotedblright{} on the ancient root \emph{h\textsubscript{1}rewd\textsuperscript{h}-} (\emph{rouge}, \emph{rosso}, Spanish \emph{rojo}, and English \emph{red}, German \emph{rot}). Portuguese went its own way.

\textbf{vermelho} $\leftarrow$ Latin \emph{\textbf{vermiculus}}, \textquotedblleft{}\textbf{little worm}\textquotedblright{} — the diminutive of \emph{vermis}, \textquotedblleft{}worm.\textquotedblright{}

Why? Because the finest red dye in the ancient Mediterranean came from \textbf{kermes}, a scale insect harvested off oak trees and crushed for its pigment. Dyers called the creature a \textquotedblleft{}little worm,\textquotedblright{} and eventually the \textbf{dye's name became the colour's name}.

The same word gave English \textbf{vermilion}. And \emph{vermis} gave English \textbf{vermin}.

Portuguese does still keep the old root — in \textbf{rubro} (literary \textquotedblleft{}red\textquotedblright{}) and \textbf{ruivo} (\textquotedblleft{}red-haired\textquotedblright{}) — but the everyday word for red is a bug.
\end{cousinweb}

\begin{cousinweb}
\textbf{azul} $\leftarrow$ Arabic \emph{\textbf{lāzaward}} $\leftarrow$ Persian \emph{\textbf{lāžward}}: \textbf{lapis lazuli}, the deep-blue stone mined for millennia in Afghanistan and ground into the most expensive pigment in the medieval world.

The initial \emph{\textbf{l-}} was lost — Spanish and Portuguese ears took it for the article \emph{l'} and dropped it — leaving \emph{azul}. The same journey produced French \textbf{azur}, Italian \textbf{azzurro}, and English \textbf{azure}.

This is the \textbf{al-Andalus} thread again: eight centuries of Arabic in Iberia left Portuguese and Spanish thousands of words, and one of them is a basic colour.
\end{cousinweb}

\begin{grammarlens}[title={What you've built: The chapter in one table}]
\noindent
\begin{tabularx}{\linewidth}{@{}>{\raggedright\arraybackslash}X>{\raggedright\arraybackslash}X@{}}
\toprule
\textbf{Portuguese} & \textbf{source} \\
\midrule
preto & Latin \emph{pressus} \textquotedblleft{}pressed\textquotedblright{} \\
\textbf{branco} & \textbf{Germanic} \emph{blank} \textquotedblleft{}shining\textquotedblright{} \\
\textbf{vermelho} & Latin \emph{vermiculus} \textquotedblleft{}\textbf{little worm}\textquotedblright{} \\
\textbf{azul} & \textbf{Arabic} \emph{lāzaward} \textquotedblleft{}\textbf{lapis lazuli}\textquotedblright{} \\
\bottomrule
\end{tabularx}

Not one of Portuguese's four basic colours is a plain inherited Latin colour word.
\end{grammarlens}

\subsection*{Your turn}
Say these aloud:

\begin{itemize}
\raggedright
  \item \textquotedblleft{}vermelho, azul\textquotedblright{}
  \item \textquotedblleft{}o vinho é vermelho… o vinho tinto\textquotedblright{} — note: \textbf{wine} is \emph{tinto}, \textquotedblleft{}dyed\textquotedblright{}
  \item \textquotedblleft{}vermelho — vermilion — vermin\textquotedblright{}
  \item all four — \textquotedblleft{}preto, branco, vermelho, azul\textquotedblright{}
\end{itemize}

\begin{tcolorbox}[breakable,colback=teal!4,colframe=teal!35!black,title={Before you move on}]
What does \textbf{vermelho} literally come from? (\emph{\textbf{Vermiculus}}, \textquotedblleft{}\textbf{little worm}\textquotedblright{} — the \textbf{kermes} insect crushed for dye.) Which English word is its sibling? (\textbf{Vermilion}.) Where does \textbf{azul} come from? (\textbf{Arabic \emph{lāzaward}}, \textquotedblleft{}\textbf{lapis lazuli},\textquotedblright{} via Persian — the \emph{l-} lost as if an article.) What's odd about all four Portuguese colours? (\textbf{None} is a plain inherited Latin colour word.) Next chapter: describing what you see.
\end{tcolorbox}
